% !TEX root = ../matan.tex

\section{Интегральные неравенства: неравенство Иенсена}

\begin{Definition}{Выпуклая (вниз) функция}{}
	Функция $\vphi:\RR\to\RR$ называется \textbf{выпуклой вниз}, если для любых
	$x,y\in\RR$ и любого $\lambda\in[0,1]$
	\[
	\vphi\bigl(\lambda x+(1-\lambda)y\bigr)\le \lambda\,\vphi(x)+(1-\lambda)\,\vphi(y).
	\]
	Геометрически: график функции лежит не выше любой своей хорды.
\end{Definition}

Нам понадобится дискретный вариант неравенства, который выражает то же
свойство для произвольной выпуклой комбинации точек.

\begin{Lemma}{Дискретное неравенство Иенсена}{}
	Пусть $\vphi:\RR\to\RR$ выпукла вниз, $\alpha_1,\dots,\alpha_N\ge0$ и
	$\sum_{k=1}^N\alpha_k=1$. Тогда для любых $x_1,\dots,x_N\in\RR$
	\[
	\vphi\!\left(\sum_{k=1}^N \alpha_k x_k\right)\le \sum_{k=1}^N \alpha_k\,\vphi(x_k).
	\]
\end{Lemma}

\begin{proof}
	Индукция по $N$. При $N=1$ имеем $\alpha_1=1$ и неравенство обращается в
	равенство. При $N=2$ это в точности определение выпуклости.

	Пусть утверждение верно для $N$ слагаемых; докажем для $N+1$. Пусть
	$\alpha_1,\dots,\alpha_{N+1}\ge0$, $\sum_{k=1}^{N+1}\alpha_k=1$. Если
	$\alpha_{N+1}=1$, то остальные нули и неравенство тривиально. Иначе положим
	$s:=\sum_{k=1}^{N}\alpha_k=1-\alpha_{N+1}\in(0,1]$ и $\beta_k:=\alpha_k/s$;
	тогда $\beta_k\ge0$ и $\sum_{k=1}^N\beta_k=1$. Обозначим
	$y:=\sum_{k=1}^N\beta_k x_k$. По определению выпуклости (для двух точек $y$ и
	$x_{N+1}$ с весами $s$ и $\alpha_{N+1}$, $s+\alpha_{N+1}=1$):
	\[
	\vphi\!\left(\sum_{k=1}^{N+1}\alpha_k x_k\right)
	=\vphi\bigl(s\,y+\alpha_{N+1}x_{N+1}\bigr)
	\le s\,\vphi(y)+\alpha_{N+1}\vphi(x_{N+1}).
	\]
	По предположению индукции $\vphi(y)\le\sum_{k=1}^N\beta_k\vphi(x_k)$, поэтому
	\[
	\vphi\!\left(\sum_{k=1}^{N+1}\alpha_k x_k\right)
	\le s\sum_{k=1}^N\beta_k\vphi(x_k)+\alpha_{N+1}\vphi(x_{N+1})
	=\sum_{k=1}^{N+1}\alpha_k\vphi(x_k),
	\]
	так как $s\beta_k=\alpha_k$. Шаг индукции завершён.
\end{proof}

\begin{Theorem}{Неравенство Иенсена (интегральная форма)}{}
	Пусть $f,g$ интегрируемы по Риману на каждом $[c,d]\subset\langle a,b\rangle$,
	причём соответствующие несобственные интегралы сходятся, и пусть
	\[
	g(x)\ge0,\qquad \int_a^b g(x)\,dx=1
	\]
	Пусть $\vphi:\RR\to\RR$ выпукла вниз.
	Тогда
	\[
	\vphi\!\left(\int_a^b g(x)f(x)\,dx\right)
	\le
	\int_a^b \vphi\bigl(f(x)\bigr)\,g(x)\,dx.
	\]
	Иными словами, $\vphi$ от среднего не превосходит среднего от $\vphi$.

	В частности, если $\langle a,b\rangle\subset\RR$ конечен, то, взяв
	$g\equiv\frac{1}{b-a}$, получаем
	\[
	\vphi\!\left(\frac{1}{b-a}\int_a^b f(x)\,dx\right)
	\le
	\frac{1}{b-a}\int_a^b \vphi\bigl(f(x)\bigr)\,dx.
	\]
\end{Theorem}

\begin{proof}
	Воспользуемся ключевым свойством выпуклых функций: \emph{в каждой точке у
	выпуклой функции есть опорная прямая}. Точнее, для выпуклой
	$\vphi:\RR\to\RR$ в любой точке $c\in\RR$ существуют конечные односторонние
	производные $\vphi'_-(c)\le\vphi'_+(c)$, и для любого
	$m\in[\vphi'_-(c),\vphi'_+(c)]$ выполнено
	\[
	\vphi(y)\ge \vphi(c)+m\,(y-c)\qquad\forall\,y\in\RR. \tag{$\ast$}
	\]
	(Это прямое следствие выпуклости: разностное отношение
	$\frac{\vphi(y)-\vphi(c)}{y-c}$ монотонно по $y$, поэтому для $y>c$ оно не
	меньше $\vphi'_+(c)$, а для $y<c$ не больше $\vphi'_-(c)$; в обоих случаях
	$(\ast)$ выполнено.)

	Положим
	\[
	c:=\int_a^b g(x)f(x)\,dx
	\]
	 и выберем $m$ из
	неравенства $(\ast)$ для этой точки $c$. Подставляя $y=f(x)$ в $(\ast)$,
	имеем для всех $x$:
	\[
	\vphi\bigl(f(x)\bigr)\ge \vphi(c)+m\bigl(f(x)-c\bigr).
	\]
	Умножим обе части на $g(x)\ge0$ и
	проинтегрируем по $\langle a,b\rangle$:
	\[
	\int_a^b \vphi\bigl(f(x)\bigr)g(x)\,dx
	\ge
	\vphi(c)\int_a^b g(x)\,dx + m\left(\int_a^b f(x)g(x)\,dx - c\int_a^b g(x)\,dx\right).
	\]
	Учитывая $\int_a^b g=1$ и $\int_a^b fg=c$, правая часть равна
	\[
	\vphi(c)\cdot1 + m\,(c-c\cdot1)=\vphi(c).
	\]
	Следовательно,
	\[
	\int_a^b \vphi\bigl(f(x)\bigr)g(x)\,dx\ge \vphi(c)
	=\vphi\!\left(\int_a^b g(x)f(x)\,dx\right),
	\]
	что и требовалось. 
\end{proof}

\begin{Remark}{}{}
	Дискретное неравенство из леммы --- частный случай интегрального: взяв
	ступенчатые $f,g$ (т.е. сосредоточив "массу" $g$ в точках с весами
	$\alpha_k$), из интегральной формы получаем
	$\vphi\!\bigl(\sum\alpha_k x_k\bigr)\le\sum\alpha_k\vphi(x_k)$.
\end{Remark}

\begin{Example}{}{}
	Для $\vphi(x)=e^{x}$ (выпукла вниз) и конечного отрезка $[a,b]$ неравенство
	Иенсена даёт неравенство между средним геометрическим и средним
	арифметическим в интегральной форме:
	\[
	\exp\!\left(\frac{1}{b-a}\int_a^b f\,dx\right)
	\le
	\frac{1}{b-a}\int_a^b e^{f(x)}\,dx.
	\]
\end{Example}
